% sections/12-conclusion.tex — §12: return to I1–I3, established vs unresolved.
\section{Conclusion}
\label{sec:conclusion}

This systematization asked whether LLM-driven vulnerability work between January 2024 and June
2026 is best understood through model capability or through something else. Our answer, built
on a 36-record multi-stream corpus of which 14 records bear autonomy classifications, is that
the field's evaluations, incentives, and discourse now measure progress in \emph{agentic
autonomy} --- while whether underlying capability moved for reasons other than scale remains,
honestly, an open question.

\textbf{On I1 (established in framing; unresolved in substance).} The autonomy gradient A0--A3
organizes the corpus cleanly: capability probes (A0) persist only as benchmark subjects;
pipeline automation (A1) runs at ecosystem scale with humans at review boundaries; task agents
(A2) appear across seeded production research, entry-point-assisted auditing, platform-operated
pentesting, and competition CRSs integrating discovery, proof, and patching. What the record
does not contain --- under our four-verb, environment-E2 operational definition --- is a single
A3 autonomous production hunter. We state the finding only as: no included evidence met our
operational definition of A3. The nearest candidate is an offensive incident,
vendor-reported with hallucination caveats and human-selected targets.

\textbf{On I2 (documented edges; causality bounded).} Offense and defense are coupled where the
record says so: Project Zero named AIxCC's Team Atlanta as the inspiration for its SQLite
variant-analysis campaign; a maintainer extended an AI-found bug's fix to sibling code paths;
\$30M-class prizes became transition funding, university labs, and foundation governance; and
threat intelligence seeded a pre-exploitation disclosure. These documented edges do not license
the arms-race vernacular: shared-cause explanations survive globally, and attacker-side
response remains unmeasured.

\textbf{On I3 (measured, increasingly binding).} Validity failures here are not rhetorical.
Information conditions swing headline exploitation rates by twelve-fold; dataset
reconstruction collapses detection SOTA; oracle gaming has a named taxonomy; environment
heuristics decide winners; and leaderboard counting conflates outcome categories. Because
these measures allocate prizes and rank vendors directly --- and increasingly inform funding
decisions through transition programs --- evaluation validity is an increasingly binding
constraint on what this field can claim to know.

\textbf{What we ask of the next iteration of this literature.} Publish information conditions
with every capability number; adopt category-aware outcome accounting; evaluate patching
against exploit-aware oracles; measure B4 rather than narrate it; and treat autonomy claims by
the weakest missing condition, not the strongest available adjective. The tools described in
this paper will be judged --- fairly or not --- by whether they make those practices easier.
